% Part: incompleteness
% Chapter: introduction
% Section: overview

\documentclass[../../../include/open-logic-section]{subfiles}

\begin{document}

\olfileid{inc}{int}{ovr}

\olsection{مروری بر نتایج ناتمامیت}

هیلبرت انتظار داشت بتوان ریاضیات را در نظریه‌ای !!{axiomatizable}
صوری کرد و سپس !!{complete} و !!{decidable}بودن آن را ثابت کرد. افزون
بر این، هدف او اثبات سازگاری این نظریه با ابزارهایی بسیار ضعیف و
«متناهی» بود تا از ریاضیات کلاسیک در برابر چالش‌های شهودگرایی دفاع
کند. قضایای ناتمامیت گودل نشان دادند این اهداف دست‌یافتنی نیستند.

قضیهٔ نخست ناتمامیت گودل نشان داد نسخه‌ای از \emph{اصول ریاضیات}
راسل و وایتهد !!{complete} نیست. اما اثبات در واقع بسیار کلی بود و بر
طیف گسترده‌ای از نظریه‌ها اعمال می‌شد. پس مشکل فقط این نبود که
\emph{اصول ریاضیات} نتوانسته است ریاضیات را کاملاً در بر گیرد؛ بلکه
\emph{هیچ} نظریهٔ پذیرفتنی چنین نمی‌کند. جداکردن ویژگی‌هایی از
نظریه‌ها که برای اعمال قضایای ناتمامیت کافی‌اند و تعمیم اثبات گودل
به‌گونه‌ای که فقط به این ویژگی‌ها وابسته باشد، زمان برد. اما اکنون
می‌توانیم صورتی بسیار کلی از قضیهٔ نخست ناتمامیت را برای نظریه‌های
زبان حساب~$\Lang L_A$ بیان کنیم.

\begin{thm}
اگر $\Gamma$ نظریه‌ای سازگار و !!{axiomatizable} در~$\Lang L_A$
باشد که همهٔ توابع محاسبه‌پذیر را !!{represents} و برای همهٔ روابط
!!{decidable} نیز چنین کند، آنگاه $\Gamma$ !!{complete} نیست.
\end{thm}

اینکه $\Gamma$ !!{complete} نیست یعنی دست‌کم یک !!{sentence}ٔ~$!A$ وجود دارد
که $\Gamma\Proves/ !A$ و~$\Gamma\Proves/\lnot !A$. چنین !!a{sentence}ای
\emph{مستقل} (از~$\Gamma$) نامیده می‌شود. در واقع می‌توانیم نسبتاً
سریع ثابت کنیم جمله‌های مستقلی باید وجود داشته باشند. اما قوت اثبات
گودل در این است که نمونه‌ای \emph{مشخص} از چنین !!{sentence}ای می‌سازد:
!!a{sentence}ای چون~$!G_\Gamma$، موسوم به \emph{جملهٔ گودلِ}~$\Gamma$، که در~$\Gamma$ با
ادعای اثبات‌ناپذیریِ خود هم‌ارز است. سازوکار معمول گودل از سازگاریِ
$\Gamma$ نتیجه می‌گیرد که $!G_\Gamma$ اثبات‌پذیر نیست؛ اما برای
نتیجه‌گیریِ اثبات‌ناپذیریِ~$\lnot !G_\Gamma$ از همان استدلال معمول،
فرضی قوی‌تر مانند $1$-سازگاری (و در صورت‌بندی اصلی،
$\omega$-سازگاری) لازم است. پالایش راسر جملهٔ مشخص دیگری می‌سازد که
استقلالش تنها از سازگاری نتیجه می‌شود و بنابراین صورت سازگاری‌محورِ
قضیهٔ بالا را به دست می‌دهد. هر دو ساخت سازنده‌اند: با داشتن
اصل‌موضوعی‌سازیِ~$\Gamma$ و توصیف دستگاه !!{derivation}، اثبات روشی برای
نوشتن واقعیِ جملهٔ متناظر ارائه می‌کند.

ساختِ اثبات گودل مستلزم آن است که راهی برای بیان ویژگی‌ها و عملیات
روی ترم‌ها و !!{formula}sی خود~$\Lang L_A$ در~$\Lang L_A$ بیابیم.
این‌ها ویژگی‌هایی مانند «$!A$ یک !!a{sentence} است» و «$\delta$ یک
!!a{derivation} از~$!A$ است»، و عملیاتی مانند~$\Subst{!A}{t}{x}$ را شامل
می‌شوند. این راه باید (الف) با «رمزگذاری» نمادها و دنباله‌هایشان (که
ترم‌ها، !!{formula}s، !!{derivation}s و مانند آن هستند) به‌صورت اعداد
طبیعی (که~$\Lang L_A$ می‌تواند درباره‌شان سخن بگوید) این ویژگی‌ها و
روابط را بیان کند؛ و (ب) این کار را چنان انجام دهد که $\Gamma$
واقعیت‌های مربوط را ثابت کند. پس باید نشان دهیم این ویژگی‌ها با
ویژگی‌های !!{decidable} اعداد طبیعی رمزگذاری می‌شوند و عملیات با
توابع محاسبه‌پذیر روی اعداد طبیعی متناظرند. این کار «حسابی‌سازی نحو»
نام دارد.

بااین‌حال، پیش از بررسی چگونگی حسابی‌سازی نحو، شرط «به‌اندازهٔ کافی
قوی» بودن~$\Gamma$، یعنی !!{represents} همهٔ توابع محاسبه‌پذیر و روابط
!!{decidable}، را بررسی می‌کنیم. این مستلزم ارائهٔ تعریفی دقیق از
«محاسبه‌پذیر» است. چنین تعریفی را به چند شیوه می‌توان داد؛ مثلاً با
مدل ماشین‌های تورینگ یا با توابعی که برنامه‌هایی در زبانی
همه‌منظوره محاسبه می‌کنند. اما چون هدف ما !!{represents}s این توابع و
روابط در نظریه‌ای به زبان~$\Lang L_A$ است، بهتر است تعریفی ساده از
محاسبه‌پذیریِ صرفاً توابع عددی برگزینیم. این همان مفهوم \emph{تابع
بازگشتی} است. پس نخست توابع بازگشتی را بررسی می‌کنیم. سپس نشان
می‌دهیم خود~$\Th{Q}$ همهٔ توابع و روابط بازگشتی را !!{represents}. بدین ترتیب می‌توانیم قضیهٔ ناتمامیت را بر نظریه‌هایی مشخص
مانند~$\Th{Q}$ و~$\Th{PA}$ اعمال کنیم، زیرا ثابت کرده‌ایم این‌ها
نمونه‌هایی از نظریه‌های «به‌اندازهٔ کافی قوی» هستند.

نتیجهٔ نهایی حسابی‌سازی نحو !!a{formula}ی چون~$\OProv[\Gamma](x)$ است که
از راه رمزگذاری !!{formula}s به‌صورت اعداد، اثبات‌پذیری از اصول
موضوعهٔ~$\Gamma$ را بیان می‌کند. مشخصاً، اگر $!A$ با عدد~$n$ رمزگذاری
شده باشد و~$\Gamma\Proves !A$، آنگاه
$\Gamma\Proves\Prov[\Gamma](\num{n})$. این «محمول اثبات‌پذیری» برای
$\Gamma$ همچنین امکان می‌دهد سازگاری~$\Gamma$ را در معنایی معین با
!!a{sentence}ای از~$\Lang L_A$ بیان کنیم: «!!{sentence}ٔ سازگاری»~$\Gamma$ را
$\lnot\Prov[\Gamma](\num{n})$ بگیرید، که در آن~$n$ رمز یک تناقض،
مثلاً~$\lfalse$، است. قضیهٔ دوم ناتمامیت می‌گوید نظریه‌های سازگار و
!!{axiomatizable} نیز جمله‌های سازگاریِ خود را ثابت نمی‌کنند. شرایط
لازم برای اعمال این قضیه اندکی سخت‌تر از صرفِ بازنمایی همهٔ توابع
محاسبه‌پذیر و روابط !!{decidable} است، اما نشان خواهیم داد
$\Th{PA}$ آن‌ها را برآورده می‌کند.

\end{document}
